\documentclass[10pt,a4paper]{amsart}
\usepackage[margin=19mm]{geometry}
\usepackage{amsmath,amssymb,mathtools}
\usepackage[scr=rsfs,cal=euler]{mathalfa}
\usepackage{xfrac}
\usepackage[unicode,hidelinks,bookmarks=false]{hyperref}
\newcommand{\ie}{i.e.\ }
\newcommand{\cf}{cf.\ }
\newcommand{\resp}{respectively\ }
\newcommand{\Subsection}{\subsection}
\providecommand{\nobreakdash}{\nobreak\hbox{-}}
\setlength{\emergencystretch}{2em}
\multlinegap0pt
\usepackage{bm}
\usepackage{bbm}
\usepackage{dsfont}
\usepackage{xspace}
\usepackage{cancel}
\usepackage{mathtools}
\usepackage{amsthm}
\makeatletter
\@ifundefined{defi}{\newtheorem{defi}{Definition}}{}
\@ifundefined{theorem}{\newtheorem{theorem}{Theorem}}{}
\makeatother
\makeatletter
\expandafter\ifx\csname add\endcsname\relax
\newenvironment{add}{\par\color{blue}}{\par}
\fi
\expandafter\ifx\csname answer\endcsname\relax
\newenvironment{answer}[1][Answer]
{\tcolorbox[noparskip,colframe=Pink,title=#1,code={\singlespacing}]}{\endtcolorbox}
\fi
\expandafter\ifx\csname question\endcsname\relax
\newenvironment{question}[1][Remark]
\fi
\expandafter\ifx\csname todiscuss\endcsname\relax
\newenvironment{todiscuss}{\par\color{red}}{\par}
\fi
\makeatother
\title[A structure-preserving multiscale solver for particle-wave i]{A structure-preserving multiscale solver for particle-wave interaction in non-uniform magnetized plasmas}
\author{Kun Huang, Irene M. Gamba, Chi-Wang Shu}
\date{Source: arXiv:2505.20210v1 (CC BY 4.0)}
\begin{document}
\maketitle

\noindent\emph{Source and licence.} A structure-preserving multiscale solver for particle-wave interaction in non-uniform magnetized plasmas. Version arXiv:2505.20210v1 (CC BY 4.0), distributed under the Creative Commons Attribution 4.0 International licence, as recorded in the arXiv licence metadata for this exact version and shown on the abstract page https://arxiv.org/abs/2505.20210v1. Full rights notice: licenses/arxiv-2505.20210-CC-BY-4.0.txt.\par\smallskip

\noindent\textit{Editorial scope.} Counted unit: the Theorem whose source label is \texttt{None} in the article (source0.tex lines 429--434, source label \texttt{None}), delivered together with the article's own complete original proof of it (source0.tex lines 436--474, one \texttt{proof} environment of 39 lines). No mathematics is written, repaired or re-derived: statement and proof are the article's own text line for line. Required context retained uncounted: connectionmat (lines 417--425). Every definition, display and label of those lines is retained unchanged. Omitted from this package and not redistributed: 1--416, 426--428, 435--435, 475--1304. Nothing inside a retained excerpt is omitted or altered: every retained body is delivered exactly as the article's lines are written, so no omission receipt of any kind accompanies this package. Citations: the retained excerpts carry no citation command at all, so no bibliography entry of the article is transcribed and no reference is redistributed. Reference adaptation: none was needed and none was made; every reference inside the delivered unit resolves inside it, so no unresolved reference or citation is produced. Printed slips kept literal: no printed word, symbol, hypothesis or number of the counted statement or of its proof was altered. Layout and class: the article's own class is \texttt{article} with packages including graphicx, inputenc, geometry, color, enumitem, authblk, lineno, hyperref, hypcap, algorithm2e, algorithmic, float, caption, subcaption; the wrapper is the lane's pinned \texttt{amsart} editorial wrapper at 10pt on A4, which restates the theorem-like environments the retained text needs and adds the article's own macro definitions that the retained text needs (none), with extra wrapper packages (bm, bbm, dsfont, xspace, cancel, mathtools). The delivered numbering is the unit's own. Assets: no retained span contains a figure environment, a graphics-inclusion command, a PDF-page inclusion, a table or a TikZ picture; no image file is copied into this package, the package declares no asset dependency and no image file is redistributed with it. Licence: version arXiv:2505.20210v1 of this article is distributed under the Creative Commons Attribution 4.0 International licence, as recorded in the arXiv licence metadata for this exact version and shown on the abstract page https://arxiv.org/abs/2505.20210v1. A full rights notice is delivered at licenses/arxiv-2505.20210-CC-BY-4.0.txt. Characters: every retained excerpt is pure ASCII, so the delivered engine, XeLaTeX with its default Latin Modern text font, reports no missing character.
\begin{defi}[connection matrix]
    For a given interval $I = (H_{a}, H_{b})$, the connection matrix associated with $I$ is defined as 
    \begin{equation}\label{connectionmat}
        A_{ij}(I)=\begin{cases}
1, & \overline{V_{i}}\cap\overline{V_{j}}\cap H^{-1}(I)\neq\emptyset\\
0, & \text{otherwise}
\end{cases}
    \end{equation}
\end{defi}

\begin{theorem}
    Suppose that $\mathcal{F}(I)$ is the family of trajectory bundles generated by interval $I=(H_{a}, H_{b})$ and $\{\mathcal{R}_{\nu} \}$ is the set of connected components determined by the connection matrix $A_{ij}(I)$ as defined in Equation(\ref{connectionmat}), then
    \begin{equation*}
        \{\mathcal{R}_{\nu}\}=\left\{ \mathcal{R}_{min}(S):S\in\mathcal{F}(I)\right\}. 
    \end{equation*}
\end{theorem}

\begin{proof}
    It takes two steps to prove the theorem.
    \begin{enumerate}
        \item Prove that $\{\mathcal{R}_{\nu}\} \subset \left\{ \mathcal{R}_{min}(S):S\in\mathcal{F}(I)\right\} $.

        By definition, any connected component $\mathcal{R}_{\nu}$ contains at least two elements, say $V_{m}$ and $V_{n}$. 
        
        Since $H^{-1}(I)$ is open,
        \begin{equation*}
            \overline{V_{m}}\cap\overline{V_{n}}\cap H^{-1}(I)\neq\emptyset\Rightarrow V_{m}\cap H^{-1}(I)\neq\emptyset.
        \end{equation*}

        Considering that
        \begin{equation*}
            \emptyset \neq V_{m}\cap H^{-1}(I)=V_{m}\cap\left(\cup_{\mathcal{F}}S\right)=\cup_{\mathcal{F}}\left(V_{m}\cap S\right),
        \end{equation*}
        
        there must exists a trajectory bundle $S\in \mathcal{F}$ such that
        \begin{equation*}
            V_{m}\cap S \neq \emptyset.
        \end{equation*}

        It remains to show that $\mathcal{R}_{\nu}=\mathcal{R}_{min}(S)$.
        \begin{itemize}
            \item Prove that $\mathcal{R}_{\nu} \subset \mathcal{R}_{min}(S)$.

            Recall that $V_{m}\cap S \neq \emptyset$, then for any $V_{l}\in \mathcal{R}_{\nu}$, there must be a continuous path inside $H^{-1}(I)$ from $(k_{1}, x_{1}) \in V_{l}\cap H^{-1}(I)$ to  $(k_{2}, x_{2}) \in V_{m}\cap H^{-1}(I)$. 
            Since $(k_{2}, x_{2}) \in S$ implies that $(k_{1}, x_{1}) \in S$, it follows that $V_{l}\cap S \neq \emptyset$. Therefore $\mathcal{R}_{\nu} \subset \mathcal{R}_{min}(S)$.
            
            \item Prove that $\mathcal{R}_{\nu} \supset \mathcal{R}_{min}(S)$.

            Consider a triangle $V_{l} \in \mathcal{R}_{min}(S)$, then by definition $V_{l} \cap S \neq \emptyset$. Suppose that $V_{l} \notin \mathcal{R}_{\nu}$, then $(k_{1}, x_{1})\in V_{l} \cap S$ is not connected to any point $(k_{2}, x_{2}) \in S\cap\cup_{i\in\mathcal{R}_{\nu}}\overline{V_{i}}$, hence $S$ is not a connected set, contradictory to its definition, therefore $V_{n} \in \mathcal{R}_{\nu}$.
        \end{itemize}
        
        \item Prove that $\{\mathcal{R}_{\nu}\} \supset \left\{ \mathcal{R}_{min}(S):S\in\mathcal{F}(I)\right\} $.

        By definition $\cup_{\nu} \mathcal{R}_{\nu} = \cup_{\mathcal{F}} \mathcal{R}_{min}(S)$, hence for any $S$ there exists an $\mathcal{R}_{\nu}$ such that $\mathcal{R}_{\nu} \cap \mathcal{R}_{min}(S) \neq \emptyset$. In the same way as above, it can be proved that $\mathcal{R}_{\nu}=\mathcal{R}_{min}(S)$.
    \end{enumerate}
\end{proof}

\end{document}
